\documentclass[10pt,a4paper]{amsart}
\usepackage[margin=19mm]{geometry}
\usepackage{amsmath,amssymb,mathtools}
\usepackage[scr=rsfs,cal=euler]{mathalfa}
\usepackage{xfrac}
\usepackage[unicode,hidelinks,bookmarks=false]{hyperref}
\newcommand{\ie}{i.e.\ }
\newcommand{\cf}{cf.\ }
\newcommand{\resp}{respectively\ }
\newcommand{\Subsection}{\subsection}
\providecommand{\nobreakdash}{\nobreak\hbox{-}}
\setlength{\emergencystretch}{2em}
\multlinegap0pt
\usepackage{bm}
\usepackage{bbm}
\usepackage{dsfont}
\usepackage{xspace}
\usepackage{cancel}
\usepackage{mathtools}
\usepackage{amsthm}
\makeatletter
\@ifundefined{theorem}{\newtheorem{theorem}{Theorem}}{}
\@ifundefined{lemma}{\newtheorem{lemma}[theorem]{Lemma}}{}
\makeatother
\title[Characterisation of Stability and Decay Rates in a Weakly Da]{Characterisation of Stability and Decay Rates in a Weakly Damped Second Order Linear Differential Equation}
\author{John A. D. Appleby, Subham Pal}
\date{Source: arXiv:2603.25547v1 (CC BY 4.0)}
\begin{document}
\maketitle

\noindent\emph{Source and licence.} Characterisation of Stability and Decay Rates in a Weakly Damped Second Order Linear Differential Equation. Version arXiv:2603.25547v1 (CC BY 4.0), distributed under the Creative Commons Attribution 4.0 International licence, as recorded in the arXiv licence metadata for this exact version and shown on the abstract page https://arxiv.org/abs/2603.25547v1. Full rights notice: licenses/arxiv-2603.25547-CC-BY-4.0.txt.\par\smallskip

\noindent\textit{Editorial scope.} Counted unit: the Theorem whose source label is \texttt{thm:asymp\_implies\_y2} in the article (source0.tex lines 813--818, source label \texttt{thm:asymp\_implies\_y2}), delivered together with the article's own complete original proof of it (source0.tex lines 820--899, one \texttt{proof} environment of 80 lines). No mathematics is written, repaired or re-derived: statement and proof are the article's own text line for line. Required context retained uncounted: lem:subexp\_conv (lines 795--798). Every definition, display and label of those lines is retained unchanged. Omitted from this package and not redistributed: 1--794, 799--812, 819--819, 900--1318. Nothing inside a retained excerpt is omitted or altered: every retained body is delivered exactly as the article's lines are written, so no omission receipt of any kind accompanies this package. Citations: the retained excerpts carry no citation command at all, so no bibliography entry of the article is transcribed and no reference is redistributed. Reference adaptation: none was needed and none was made; every reference inside the delivered unit resolves inside it, so no unresolved reference or citation is produced. Printed slips kept literal: no printed word, symbol, hypothesis or number of the counted statement or of its proof was altered. Layout and class: the article's own class is \texttt{amsart} with packages including amsmath, amssymb, amsthm, geometry; the wrapper is the lane's pinned \texttt{amsart} editorial wrapper at 10pt on A4, which restates the theorem-like environments the retained text needs and adds the article's own macro definitions that the retained text needs (none), with extra wrapper packages (bm, bbm, dsfont, xspace, cancel, mathtools). The delivered numbering is the unit's own. Assets: no retained span contains a figure environment, a graphics-inclusion command, a PDF-page inclusion, a table or a TikZ picture; no image file is copied into this package, the package declares no asset dependency and no image file is redistributed with it. Licence: version arXiv:2603.25547v1 of this article is distributed under the Creative Commons Attribution 4.0 International licence, as recorded in the arXiv licence metadata for this exact version and shown on the abstract page https://arxiv.org/abs/2603.25547v1. A full rights notice is delivered at licenses/arxiv-2603.25547-CC-BY-4.0.txt. Characters: every retained excerpt is pure ASCII, so the delivered engine, XeLaTeX with its default Latin Modern text font, reports no missing character.
	\begin{lemma}[Convolution with Subexponential Weights] \label{lem:subexp_conv}
		Let $\rho$ be a subexponential function. Let $k$ be such that there exist constants $C > 0$ and $\alpha > 0$ such that $|k(t)| \le C e^{-\alpha t}$ for all $t \ge 0$. 
		If a continuous function $z$ satisfies $z(t) = O(\rho(t))$ as $t \to \infty$, then $(k * z)(t) = O(\rho(t))$ as $t \to \infty$.
	\end{lemma}

	\begin{theorem}[Asymptotic Form Implies $o(1/A(t))$ Residual] \label{thm:asymp_implies_y2}
		Let $\rho(t) = 1/A(t)$. Assume $p(t) \to 0$ and $p'(t) \to 0$ as $t \to \infty$. 
		If the physical state satisfies
		\[ x(t) = x_0(t) + x_1(t) = \rho(t)\left[c_1 \sin(\omega t) + c_2 \cos(\omega t)\right] + o(\rho(t)) \quad \text{as } t \to \infty, \]
		where $x_0(t)$ is the leading oscillatory term and $x_1(t) = o(\rho(t))$ is the remainder, then the second-order filtered forcing satisfies $y_2(t) = o(\rho(t))$.
	\end{theorem}

	\begin{proof}
		Recall the representation of $y_2(t)$ in terms of $x(t)$ using the exponentially decaying kernel $h(t)$ and its derivatives:
		\[ y_2(t) = x(t) + (h'' * x)(t) + \omega^2(h * x)(t) + (h' * px)(t) - (h * p'x)(t). \]
		By hypothesis, $h$, $h'$, and $h''$ are all bounded by $C e^{-\alpha t}$ for some constants $C, \alpha > 0$. We will analyze the terms of $y_2(t)$ in three distinct parts.
		
		Because $x_1(t) = o(\rho(t))$, Lemma \ref{lem:subexp_conv} immediately implies that convolutions of $x_1$ with any exponentially decaying kernel remain $o(\rho(t))$. Thus, $x_1 + h'' * x_1 + \omega^2 h * x_1 = o(\rho(t))$.
		
		Next, consider the terms involving $p(t)$ and $p'(t)$. We know $x(t) = O(\rho(t))$. Since $p(t) \to 0$ and $p'(t) \to 0$, the products are strictly smaller than $\rho(t)$:
		\[ p(t)x(t) = o(\rho(t)) \quad \text{and} \quad p'(t)x(t) = o(\rho(t)). \]
		Applying Lemma \ref{lem:subexp_conv} again, the convolutions $(h' * px)(t)$ and $(h * p'x)(t)$ are strictly $o(\rho(t))$. 
		
		Therefore, all that remains is to evaluate the linear leading-order operator $L$ applied to the main oscillatory term $x_0(t)$:
		\[ L[x_0](t) = x_0(t) + (h'' * x_0)(t) + \omega^2 (h * x_0)(t). \]	
		Now, we recall that  $\rho(t) = \exp\left(-\frac{1}{2}\int_0^t p(s)\,ds\right)$, where $\lim_{t \to \infty} p(t) = 0$ and $\lim_{t \to \infty} p'(t) = 0$, and  $x_0(t) = \rho(t)\Theta(t)$, where $\Theta(t) = c_1 \cos(\omega t) + c_2 \sin(\omega t)$. Recall also the definition 
	   $h(t) = t e^{-\omega^2 t}$ for some $\omega > 0$. We must show that 
\[ L[x_0](t)=: l(t) := x_0(t) + (h'' * x_0)(t) + \omega^2 (h * x_0)(t).
\]
is such that $l(t)=o(\rho(t))$. 
First, note that 			
$\rho'(t) = o(\rho(t))$, $\rho''(t) = o(\rho(t))$ as $t \to \infty$. By definition, we have 
\[ \rho'(t) = %\exp\left(-\frac{1}{2}\int_0^t p(s)\,ds\right) \cdot \left(-\frac{1}{2}p(t)\right) = 
-\frac{1}{2}p(t)\rho(t), \]
so, as $p(t)\to 0$ as $t\to\infty$, $\rho'(t)=o(\rho(t))$. 
Differentiating again we get 	
			\begin{align*}
				\rho''(t) &= -\frac{1}{2}p'(t)\rho(t) - \frac{1}{2}p(t)\rho'(t) = \rho(t) \left( \frac{1}{4}p(t)^2 - \frac{1}{2}p'(t) \right).
			\end{align*}
			Since $\lim_{t \to \infty} p(t) = 0$ and $\lim_{t \to \infty} p'(t) = 0$, we have  $\rho''(t) = o(\rho(t))$. 
			
Next, we analyse the $x_0$ dependent terms in $l$. 
			Notice that $\Theta(t)$ obeys $\Theta''(t) + \omega^2 \Theta(t) = 0$. Then 
			\[ x_0'(t) = \rho'(t)\Theta(t) + \rho(t)\Theta'(t), \]
			and 
			\[ x_0''(t) = \rho''(t)\Theta(t) + 2\rho'(t)\Theta'(t) + \rho(t)\Theta''(t) \]
			Substituting $\Theta''(t) = -\omega^2 \Theta(t)$ into the second derivative, we get 
			%:
			%\[ x_0''(t) = \rho''(t)\Theta(t) + %2\rho'(t)\Theta'(t) - \omega^2 x_0(t). \]
			%Rearranging this, we construct the forced harmonic expression:
			\[ x_0''(t) + \omega^2 x_0(t) = \rho''(t)\Theta(t) + 2\rho'(t)\Theta'(t). \]
			Since $\Theta(t)$ and $\Theta'(t)$ are bounded trigonometric functions, and we have proven $\rho' = o(\rho)$ and $\rho'' = o(\rho)$, it immediately follows that:
			\[ x_0''(t) + \omega^2 x_0(t) = o(\rho(t)). \]
			Let us define $g(t) := x_0''(t) + \omega^2 x_0(t)$. We have established $g(t) = o(\rho(t))$.
				
			We are given $h(t) = t e^{-\omega^2 t}$. Then $h(0)=0$, $h'(0)=1$. We now examine the convolution $(h'' * x_0)(t) = \int_0^t h''(t-s)x_0(s)\,ds$. We apply integration by parts twice to shift the derivatives from $h$ to $x_0$. 
			First integration by parts:
			\[ \int_0^t h''(t-s)x_0(s)\,ds = \Big[ -h'(t-s)x_0(s) \Big]_{s=0}^{s=t} + \int_0^t h'(t-s)x_0'(s)\,ds \]
			\[ = -h'(0)x_0(t) + h'(t)x_0(0) + \int_0^t h'(t-s)x_0'(s)\,ds. \]
			Second integration by parts on the remaining integral:
			\[ \int_0^t h'(t-s)x_0'(s)\,ds = \Big[ -h(t-s)x_0'(s) \Big]_{s=0}^{s=t} + \int_0^t h(t-s)x_0''(s)\,ds \]
			\[ = -h(0)x_0'(t) + h(t)x_0'(0) + (h * x_0'')(t). \]
			Combining these and applying the boundary values $h(0) = 0$ and $h'(0) = 1$:
			\[ (h'' * x_0)(t) = -x_0(t) + h'(t)x_0(0) + h(t)x_0'(0) + (h * x_0'')(t). \]
			Now, substitute this expanded form back into the definition of $l(t)$:
			\begin{align*}
				l(t) &= x_0(t) + \Big( -x_0(t) + h'(t)x_0(0) + h(t)x_0'(0) + (h * x_0'')(t) \Big) + \omega^2 (h * x_0)(t) \\
				&= h'(t)x_0(0) + h(t)x_0'(0)T(t) + \Big( h * (x_0'' + \omega^2 x_0) \Big)(t) \\
				&=: T(t) + (h * g)(t).
			\end{align*}
			
			We must show that both $T(t)$ and $(h * g)(t)$ are $o(\rho(t))$. 
			
			\textbf{1. The Transients:} 
			$T(t) = (1-\omega^2 t)e^{-\omega^2 t}x_0(0) + t e^{-\omega^2 t}x_0'(0)$. This decays exponentially. Since $p(t) \to 0$, the integral $\int_0^t p(s)\,ds$ grows slower than any linear function $\epsilon t$, meaning $\rho(t)$ decays at most sub-exponentially. Therefore, the exponentially decaying $T(t)$ is trivially $o(\rho(t))$.
			
			\textbf{2. The Convolution:}
			We must evaluate $\lim_{t \to \infty} \frac{1}{\rho(t)} \int_0^t h(s) g(t-s) \, ds$.
			Because $g(t) = o(\rho(t))$, for any $\epsilon > 0$, there exists an $M > 0$ such that for all $\tau > M$, $|g(\tau)| < \epsilon \rho(\tau)$. Furthermore, $g(t)$ is globally bounded by $C \rho(t)$ for some constant $C$.
			We split the integral at $t-M$:
			\[ \frac{1}{\rho(t)} \int_0^t |h(s)| |g(t-s)| \, ds = \int_0^{t-M} |h(s)| \frac{|g(t-s)|}{\rho(t)} \, ds + \int_{t-M}^t |h(s)| \frac{|g(t-s)|}{\rho(t)} \, ds. \]
			
			In the first integral ($s \in [0, t-M]$), we have $t-s \ge M$, so $|g(t-s)| < \epsilon \rho(t-s)$. Note that:
			\[ \frac{\rho(t-s)}{\rho(t)} = \exp\left( \frac{1}{2}\int_{t-s}^t p(\tau)\,d\tau \right). \]
			Since $p(t) \to 0$, $p(t)$ is bounded by some small constant $P \ll \omega^2$. Thus, the ratio $\rho(t-s)/\rho(t)$ is bounded by $e^{Ps/2}$. Because $h(s) = s e^{-\omega^2 s}$ decays much faster than $e^{Ps/2}$ grows, the product $|h(s)|\frac{\rho(t-s)}{\rho(t)}$ is entirely absolutely integrable on $[0, \infty)$. Let this integral be bounded by a constant $K$. The first integral is thus strictly bounded by $\epsilon K$.
			
			For the second integral ($s \in [t-M, t]$), $h(s)$ is evaluated at values tending to infinity. Since $h(s)$ decays exponentially, and the integration interval $[t-M, t]$ has fixed length $M$, this integral vanishes completely as $t \to \infty$.
			
			Since $\epsilon$ was arbitrary, the entire convolution is bounded by $0$ in the limit. Thus, $(h * g)(t) = o(\rho(t))$.
			
			Summing the asymptotic bounds, $l(t) = T(t) + (h * g)(t) = o(\rho(t)) + o(\rho(t)) = o(\rho(t))$, completing the proof.
	\end{proof}

\end{document}
